\documentclass[11pt]{article}
\pdfinfoomitdate=1
\pdftrailerid{}
\pdfsuppressptexinfo=-1
\usepackage[margin=1in]{geometry}
\usepackage{amsmath,amssymb,amsthm,mathtools}
\usepackage{enumitem,microtype,booktabs}
\usepackage[hidelinks]{hyperref}


\newtheorem{theorem}{Theorem}[section]
\newtheorem{lemma}[theorem]{Lemma}
\newtheorem{proposition}[theorem]{Proposition}
\newtheorem{corollary}[theorem]{Corollary}
\newtheorem{definition}[theorem]{Definition}
\newtheorem{remark}[theorem]{Remark}
\newcommand{\N}{\mathbb N_0}
\newcommand{\E}{\mathbb E}
\newcommand{\Pp}{\mathbb P}
\newcommand{\cL}{\mathcal L}
\newcommand{\cC}{\mathcal C}
\newcommand{\cR}{\mathcal R}
\newcommand{\one}{\mathbf 1}

\title{Positive Recurrence for Bimolecular Weakly Reversible\\
Stochastic Reaction Networks with Three Species\\
and Two Linkage Classes}
\author{Alec Kriebel\\\small Independent researcher\\
\small\href{mailto:me@aleckriebel.com}{me@aleckriebel.com}}
\date{Version 1.0 --- August 2026}

\begin{document}
\maketitle

\begin{abstract}
We prove positive recurrence on every closed communicating class for every
finite weakly reversible stochastic mass-action network with molecularity at
most two, at most three dynamically active species, at most two active
linkage classes, and arbitrary positive rate constants.  The proof combines
an embedded labelled reaction-count Green construction with a finite
source-rate compactification.  The new stochastic ingredient is an aggregate
physical workload-debt excursion: positive reactions are retained at their
natural probabilities, their actual targets initiate faster return
relaxations, and repeated slower arrivals are controlled by a scalar reflected
debt process rather than by genealogical credits.  A finite phase
classification supplies service, structural exit, or an affine invariant.
The one-active case is reduced to an exact polynomial-rate phase chain, while
the two-active boundary geometry is closed by an independently reproduced
finite atlas.  Embedded return, physical return, and nonexplosion are treated
separately.
\end{abstract}

\noindent\textbf{Keywords.}
stochastic reaction network; weak reversibility; positive recurrence;
mass-action kinetics; Green occupation; Foster--Lyapunov method; finite phase.

\section{Model and theorem}

Let the species set have at most three dynamically active coordinates.  A
complex is \(y\in\N^d\), \(d\le3\), and the network is
\emph{bimolecular} when \(|y|_1\le2\) for every complex.  A labelled reaction
channel is \(e:y_e\to y_e'\), with rate \(\kappa_e>0\), reaction vector
\(\zeta_e=y_e'-y_e\), and stochastic mass-action propensity
\[
 \lambda_e(x)=\kappa_e(x)_{y_e},\qquad
 (x)_y=\prod_i\frac{x_i!}{(x_i-y_i)!}
\]
when \(x\ge y\), and zero otherwise.  Parallel channels are retained.
Each linkage class is a connected component of the undirected complex graph;
weak reversibility means that each linkage class is strongly connected.

Fix one closed irreducible population class \(\Gamma\).  The population CTMC
has formal generator
\[
 \cL f(x)=\sum_e\kappa_e(x)_{y_e}
 \bigl[f(x+\zeta_e)-f(x)\bigr].
\]

\begin{theorem}[T3--2 positive recurrence]\label{thm:main}
Suppose the network has at most three dynamically active species and at most
two active linkage classes.  Then, for every positive rate vector, every
closed communicating class is positive recurrent.  Equivalently, every
closed irreducible class has a unique stationary probability distribution.
\end{theorem}

The one-linkage case is supplied by the marked-target theorem
\cite{Kriebel2026Single}; earlier binary one-linkage results under additional
structural assumptions appear in \cite{AndersonCappellettiKim2020}.  The
present proof uses that theorem only at explicit one-linkage reductions.

\section{Labelled embedded Green occupation}

Let \(Z_n\) be the population immediately after the \(n\)-th reaction and
\(E_{n+1}\) the actual labelled channel producing \(Z_{n+1}\).  Conditional on
\(Z_n=x\),
\[
 \Pp(E_{n+1}=e\mid Z_n=x)
 =\frac{\kappa_e(x)_{y_e}}{\Lambda(x)},\qquad
 \Lambda(x)=\sum_f\kappa_f(x)_{y_f}.
\]

Fix \(o\in\Gamma\), an increasing finite exhaustion \(D_M\ni o\), and
\[
 T_o^+=\inf\{n\ge1:Z_n=o\},\quad
 \sigma_M=\inf\{n\ge0:Z_n\notin D_M\},\quad
 \tau_M=T_o^+\wedge\sigma_M.
\]

\begin{theorem}[Embedded labelled Green occupation]\label{thm:green}
If \(\E_oT_o^+=\infty\), then \(t_M=\E_o\tau_M\to\infty\) and
\[
 \nu_M(x,e)=\frac1{t_M}\E_o\sum_{n<\tau_M}
 \one\{Z_n=x,E_{n+1}=e\}
\]
is a probability measure on labelled transitions.  It escapes every fixed
finite set.  For bounded \(f\),
\[
 \sum_{x,e}\nu_M(x,e)\,[f(x+\zeta_e)-f(x)]
 =\frac{\E_o[f(Z_{\tau_M})-f(o)]}{t_M}.
\]
For every linear workload \(h\),
\[
 \sum_{x,e}\nu_M(x,e)\,h\!\cdot\!\zeta_e=h\!\cdot q_M,\qquad
 q_M=\frac{\E_o[Z_{\tau_M}-o]}{t_M}.
\]
The vectors \(q_M\) are bounded and every limit belongs to
\(\mathbb R_{\ge0}^d\).
\end{theorem}

\begin{proof}
Monotone convergence gives \(t_M\uparrow\E_oT_o^+=\infty\).
For a fixed state \(x\), choose a finite positive-probability embedded path
from \(x\) to \(o\) not revisiting \(x\).  Each visit to \(x\) has a fixed
positive chance of terminating the killed excursion before the next visit,
so its expected visit count is uniformly bounded in \(M\).  Division by
\(t_M\) proves escape from finite sets.

The identities are pathwise telescoping.  Since each reaction changes each
coordinate by at most two,
\(\|Z_{\tau_M}-o\|_\infty\le2\tau_M\), so \(q_M\) is bounded.  Also
\(q_{M,i}\ge-o_i/t_M\), proving nonnegativity of every limit.
\end{proof}

\section{Terminal chart localization}

A Green occupation is a measure, not a representative sequence.  We
therefore compactify the occupation itself.

For each transition source record:
\begin{enumerate}[label=(\roman*)]
\item each population coordinate in \(\N\cup\{\infty\}\);
\item capped availability \(\min(x_i,2)\) and enabled-source support;
\item the actual current target and labelled channel;
\item the complete recursively normalized physical source-propensity flag;
\item the tied positive ratio vector, workload chamber, and lattice phase.
\end{enumerate}
The flag compactification is finite-depth: normalize the finite source vector;
then normalize recursively on its zero coordinates.  Thus every sequence has
a convergent flag subsequence, including arbitrarily slow divergent species.

\begin{lemma}[Terminal chart]\label{lem:chart}
After subsequence extraction, every escaping Green occupation has a nonzero
restriction to a terminal chart with fixed active set, fixed finite box for
inactive coordinates, fixed enabled support, fixed source flag, compact tied
ratio cell, fixed workload chamber and zero normalized structural-exit flux.
For the positive integer workload \(H(x)=h\cdot x\) of that chart,
\[
 \sum_{x,e}\nu(x,e)\,h\cdot\zeta_e\ge0.       \tag{3.1}
\]
\end{lemma}

\begin{proof}
Push \(\nu_M\) to the compact descriptor.  A weak limit exists and has total
mass one at infinity.  Bounded jumps preserve each asymptotic descriptor on
a positive-mass limiting edge.  Splitting over the finitely many active sets
prevents, for example, mass near \((n,0,0)\) and \((0,n,0)\) from being put
in one chart that enables both active families.

Inactive coordinates have tight probability marginals.  Choose a finite box
carrying all but \(\varepsilon\) of the source and endpoint mass and pad by
the maximal jump.  Box crossings, support changes and flag refinements are
retained as chart edges.  Indicator telescoping gives a finite flow.  Select
a terminal strongly connected component of its support and diagonalize as
\(\varepsilon\downarrow0\).

For (3.1), restrict to a long integer workload band.  Summing lower-cut flux
over consecutive cuts counts each bounded jump only a bounded number of
times.  Hence one lower cut has vanishing normalized inward flux.  Upper-cut
flux is outward.  The finite chart-error tends to zero with \(\varepsilon\).
\end{proof}

A raw reaction-count normalization can be dominated by faster neutral
motion.  We therefore take exact physical traces at successively slower
source layers.  If a layer has bounded expected count, its bounded
displacement cannot support escape.  If its count diverges, normalize by that
count and contract the already processed finite neutral phase.  The source
flag is finite, so this procedure terminates at the first occupied changing
trace.


\subsection{The iterated source simplex}

For completeness we spell out the finite compactification used in
Lemma~\ref{lem:chart}.  Let \(\mathcal S_0\) be the finite set of source
complexes enabled in the selected capped phase, and put
\[
 a_s(x)=\bar\kappa_s(x)_s,\qquad s\in\mathcal S_0.
\]
The first projective vector is
\[
 \rho^{(1)}_s(x)=\frac{a_s(x)}{\sum_{u\in\mathcal S_0}a_u(x)}.
\]
At a boundary limit let \(\mathcal F_1\) be its positive support.  On
\(\mathcal S_1=\mathcal S_0\setminus\mathcal F_1\), provided it is nonempty,
renormalize again:
\[
 \rho^{(2)}_s(x)=\frac{a_s(x)}{\sum_{u\in\mathcal S_1}a_u(x)}.
\]
Continue.  At most \(|\mathcal S_0|\) steps are possible.  The resulting
ordered partition
\[
 \mathcal F_1\gg\mathcal F_2\gg\cdots\gg\mathcal F_r
\]
and the positive ratio vector in every \(\mathcal F_j\) constitute the
complete flag.  This is a closed subset of a finite product of simplices and
is compact.  Sources with vanishing first-order ratio are not discarded; they
occur in a later \(\mathcal F_j\).

The mass-action source exponents belong to the finite set
\[
 \{0,e_i,2e_i,e_i+e_j\}.
\]
After passing to one flag cell, its strict comparisons can be represented by
a rational workload.  Indeed, take the finitely many rational separating
covectors produced by the successive normalized-log levels, clear
denominators, and combine them with a sufficiently large integer base.
The resulting \(h\in\mathbb Z_{\ge0}^d\) preserves every strict source
comparison and is positive on every active coordinate.  Equality of source
orders is retained because the complete flag has already separated every
nonzero lower-order difference.

\subsection{Why chart flow is retained}

Let \(\mathcal P\) be a finite partition of one compact flag cell into its
finite discrete data.  For a part \(A\in\mathcal P\), apply the bounded-test
identity of Theorem~\ref{thm:green} to \(\one_A\).  The normalized inflow
minus outflow tends to zero.  Hence every weak limit induces a finite
circulation after source and return errors vanish.  If a transition changes
the active set, bounded box, source support, flag face or workload chamber,
we place it on a directed chart edge rather than deleting it.  A terminal
strongly connected component therefore has zero limiting flux through every
such edge.

The workload-band argument is also purely reaction-count based.  If
\(J=\max_e|h\cdot\zeta_e|\), a transition crosses at most \(J\) integer lower
cuts.  Summing the normalized inward flux over \(K\) consecutive cuts is at
most \(J\).  One cut has inward flux at most \(J/K\).  Letting \(K\to\infty\)
proves the local nonnegative outward balance without integrating an
unbounded propensity.

\section{Actual targets and aggregate workload debt}

At an embedded jump, retain the target \(T_n\) of the actual labelled
reaction.  It is physically present.  The inherited residual factorial
identity is
\[
 V(x,t)=\sum_i\log((x_i-t_i)!),\qquad
 \Delta V=\log\frac{(x)_t}{(x)_s}.
\]
It is used only for same-linkage target paths beginning from an already
carried target; no future activation is conditioned upon.

For the source-layer argument, scale the chart workload so that \(H(y)\) is
integer-valued on complexes.  A positive reaction \(y\to u\) has \(H(u)>H(y)\)
and hence creates an actual target in a strictly faster source layer.

\begin{lemma}[Actual-target return prefix]\label{lem:return}
For every positive reaction \(y\to u\), a directed path in its linkage from
\(u\) back to \(y\) has a first target \(v\) with \(H(v)\le H(y)\).  Every
earlier source is strictly faster than \(y\), and
\[
 H(u-y)+H(v-u)=H(v-y)\le0.                    \tag{4.1}
\]
The lifted path is sequentially enabled until a declared chart exit.
\end{lemma}

\begin{proof}
Strong connectivity supplies the path.  The first-crossing definition gives
the priority statement and (4.1).  Each designated reaction produces the
source of the next reaction, so the lifted path is physical.
\end{proof}

The bounded coordinates, actual target, availability, source support,
linkage and lattice data form a finite phase.  Include every processed
physical reaction and every return prefix.

\begin{lemma}[Finite physical service phase]\label{lem:phase}
A closed finite phase component without a negative workload edge contains no
positive workload edge.  Consequently each component is strict, edgewise
zero, or has a structural exit.
\end{lemma}

\begin{proof}
A positive edge has the return certificate of Lemma~\ref{lem:return}.
Closedness retains that certificate.  A finite path whose initial reward is
positive and whose cumulative reward is nonpositive contains a negative edge.
\end{proof}

Finite graph reachability and compactness of the tied-rate cell give constants
\(L<\infty\), \(q>0\), and \(K<\infty\): in a strict component a negative
edge is reached within \(L\) carrier changes with probability at least \(q\);
the number of changes has geometric moments.  Reactions leaving the selected
source enabled do not delay its physical exponential clock.  If the total
unprocessed slower hazard is at most \(\varepsilon_R\) times the carrier
hazard, then
\[
 \Pp\{\text{slower interruption before service}\}\le K\varepsilon_R,
 \qquad \varepsilon_R\to0.                    \tag{4.2}
\]

Fix an excursion baseline \(B=H(X_0)\) and define only the physical scalar
\[
 D=(H(X)-B)^+.
\]

\begin{theorem}[Aggregate-debt clearing]\label{thm:debt}
At successive service trials suppose
\[
 D_{k+1}\le(D_k-S_k)^++A_k,
\]
and, on \(D_k>0\),
\[
 \Pp(S_k\ge1\mid\mathcal F_k)\ge p,\qquad
 \E(A_k\mid\mathcal F_k)\le a<p.
\]
Then
\[
 \E\tau_0\le \frac{D_0}{p-a}.
\]
If trial durations have conditional mean at most \(T\), the physical clearing
time has mean at most \(TD_0/(p-a)\).
\end{theorem}

\begin{proof}
For integer \(d>0\),
\[
 (d-s)^++a-d\le-\one_{\{s\ge1\}}+a.
\]
The conditional drift is at most \(-(p-a)\).  Sum through
\(n\wedge\tau_0\), use nonnegativity, and apply monotone convergence.
\end{proof}

Here \(A_k\) is the workload of the first unprocessed slower event.
Equation (4.2) gives \(\E A_k=O(\varepsilon_R)\); choose the chart threshold
so that it is below \(p/2\).  Multiple slower arrivals occur over repeated
trials and are all represented in \(D\).

After one service, if \(D>0\), restart from the new physical state and actual
target.  A strict phase gives the same service minorization; loss of support
is an exit; a zero phase passes the unchanged debt to the next slower physical
trace.  No creator identity is retained.

At \(D=0\), a strict negative transition before the next positive event lowers
the workload below the baseline.  A failed trial creates bounded debt, which
is cleared before the next zero-debt trial.  Geometric trials therefore give
genuine descent in finite mean.


\subsection{The physical carrier graph}

We give the construction behind the finite phase assertion.  A carrier state
records the actual target complex, the exact inactive-coordinate population,
the capped active availability, the creator linkage and a position in one of
the fixed return prefixes.  It does not record an active population count.
For each processed physical reaction there are two cases.

If the selected return source remains enabled, its exponential clock
continues.  Reactions not disabling that source are skipped in the carrier
clock.  If a reaction disables the source, then, inside a padded chart, it
must consume one of the finitely many inactive cofactors required by the
source.  Its actual target and the new inactive population determine another
carrier state.  Loss of the padded box or active threshold is an exit.

Thus the carrier state space is finite.  From every positive edge, the
designated return prefix gives a path to cancellation or exit.  Add all
possible disabling transfers.  A closed class avoiding cancellation would
contain the positive edge but, by the return certificate, no negative edge,
contradicting Lemma~\ref{lem:phase}.  This proves the qualitative
reactivation alternative directly from physical channels.

For uniformity, eliminate source priorities recursively along the same
complete flag.  A strictly faster recurrent class either already contains
service or is edgewise workload-neutral and is contracted by its finite
Green/Poisson corrector.  Consequently, when a carrier source at a given
priority is inspected, only sources tied with it remain as genuine
competitors; faster activity has already been absorbed into the finite phase.

Choose in every strict carrier class one path to service.  There are finitely
many phases and finitely many selected edges.  Conditional channel
probabilities from a fixed source depend only on its positive rate constants,
and tied-source ratios range over a compact cell.  Their minimum is a number
\(q>0\).  If the longest selected path has length \(L\), one block succeeds
with probability at least \(q^L\).  Repeated blocks give a geometric bound
on carrier changes.  This is a trace argument, not an assumption of fast
mixing.

\subsection{Slower arrivals}

Let \(K\) bound the expected number of carrier changes in one strict trial.
At every carrier phase, let \(\lambda_f\) be the total processed carrier
hazard and \(\lambda_s\) the total unprocessed source hazard.  The complete
flag gives
\[
 \lambda_s\le\varepsilon_R\lambda_f,\qquad \varepsilon_R\downarrow0.
\]
Stop at the first unprocessed event.  Compensator comparison over the
carrier occupation gives
\[
 \Pp\{\hbox{interruption}\}\le K\varepsilon_R.   \tag{4.3}
\]
If \(\Delta_+\) is the largest positive integer workload jump, then
\[
 \E A\le K\Delta_+\varepsilon_R,\qquad
 \E z^A\le1+K\varepsilon_R(z^{\Delta_+}-1),\quad z>1. \tag{4.4}
\]
There is at most one unprocessed event in a trial; multiple arrivals are
handled over successive trials.  This is why the arrival bound is independent
of the current debt.

For a finite workload band, stop also at its padded upper boundary.  One
reaction overshoots that boundary by at most \(\Delta_+\).  A partial
carrier excursion at the Green stopping time is therefore either a recorded
upper-shell exit or has a uniformly bounded workload correction.  When the
layer trace count diverges, this one terminal correction vanishes after
normalization.  Thus no unfinished busy period is silently omitted.

\subsection{Reactivation and zero debt}

The service selector is recomputed after every trial from the current
physical phase.  In a strict component the same finite minorization applies.
A zero component is carried to the next source layer, and any loss of the
phase support is an exit.  Therefore Theorem~\ref{thm:debt} can be restarted
whenever \(D>0\), without identifying which reaction created the remaining
excess.

At \(D=0\), run the same strict trial.  A strict edge before the first
unprocessed positive event gives a unit of descent below the baseline.  On a
failed trial, the first unprocessed event creates at most \(\Delta_+\) debt.
Clear it by Theorem~\ref{thm:debt}, then retry.  The number of trials is
geometrically dominated and all physical durations have finite mean.

\section{Finite source-layer hierarchy}

Process the complete source flag from fastest to slowest.

At the fastest layer, a positive edge is impossible: its target would be a
faster source.  Thus every recurrent class is strict or zero.

Inductively assume faster layers have been replaced by finite-mean effective
transitions with nonpositive complete workload.  Retain a raw event from the
next layer at its natural probability.  If it is positive, its actual target
starts a faster debt excursion.  Same-layer and slower interruptions enter
the arrival variable in Theorem~\ref{thm:debt}.  Completion gives a
nonpositive effective transition or a structural exit.

\begin{proposition}[Source-layer alternative]\label{prop:layers}
Every recurrent effective source-layer class has exactly one of:
\begin{enumerate}[label=(\roman*)]
\item a strict negative stationary corrected reward;
\item edgewise zero reward and a bounded rational phase coboundary;
\item positive structural-exit flux.
\end{enumerate}
The hierarchy terminates after finitely many layers.
\end{proposition}

\begin{proof}
All effective rewards are nonpositive.  In a finite irreducible class, one
negative edge has positive stationary flow and makes the mean strict.
If the mean is zero, every positive-probability edge has reward zero.  Cycle
sums vanish, so path sums define the phase coboundary.  A layer of vanishing
fraction is retained as the next slower trace.  There are finitely many
source layers.
\end{proof}

If every layer is zero, \(H(X)+\psi(\phi)\) is invariant on the terminal
recurrent support.  The phase term is bounded and \(h\) is positive on every
active coordinate, contradicting escape.  Because zero is edgewise after
complete faster relaxation, the corrected variance is also zero; no hidden
critical random walk remains.


\subsection{Induction inside one raw source layer}

The induction has an inner priority induction.  At the fastest source
priority, a positive reaction is impossible because its target would be a
strictly faster source.  Hence the fastest raw edges are neutral or negative.

Suppose priorities above \(p\) have already been replaced by nonpositive
finite-mean effective transitions.  A positive raw reaction with source
priority \(p\) produces an actual target of priority greater than \(p\).
Start the carrier process immediately after that actual reaction, while
retaining the reaction's positive workload in \(D\).  A second event from
priority \(p\) or below is an unprocessed arrival.  Equations
(4.3)--(4.4) and Theorem~\ref{thm:debt} show that the entire positive event
plus faster relaxation has nonpositive endpoint workload in finite mean, or
exits the chart.  It may cancel exactly; it is not called strict descent.

Descending through the finite priority set converts every raw edge of the
source layer into a nonpositive effective edge.  Only then is the finite
recurrent-class alternative of Proposition~\ref{prop:layers} applied.

\subsection{Passage to a slower reaction-count trace}

Let \(C_j(M)\) be the expected number of layer-\(j\) events in the killed
Green path.  If \(C_j(M)\) is bounded, its bounded jumps have zero
contribution to an escaping endpoint after division by the relevant divergent
count.  If \(C_j(M)\to\infty\), normalize its exact labelled trace, including
the complete faster excursions between successive layer-\(j\) events.
Positive chart-exit counts are retained.  This gives a nonzero finite-phase
occupation at layer \(j\).

Thus ``zero at the current normalization'' never means deletion.  It means
that the same physical events appear at the next slower normalization.
Since the source set is finite, the procedure cannot continue indefinitely.

\subsection{The all-zero branch}

Suppose every effective layer is edgewise zero.  On each finite recurrent
phase class, define \(\psi\) by integrating \(-h\cdot\Delta X\) along phase
paths.  Zero cycle sums make this well-defined.  Correctors from successive
layers add, and only finitely many are present.  Hence
\[
 h\cdot X+\psi(\phi)
\]
is invariant for every transition of the terminal physical trace.  Events
outside the trace have bounded count or positive exit flux.  The first cannot
produce an escaping endpoint and the second contradicts terminality.  Since
\(\psi\) is bounded, the invariant excludes escape in an active direction.

\section{Active-coordinate cases}

\subsection{Three active coordinates}

Every possible bimolecular source is enabled.  At the first occupied changing
layer, every changing reaction is nonincreasing and at least one is strict:
a positive reaction at the highest changing source level would force, along
its linkage return cycle, a negative reaction sourced at a still higher
level.  Compact tied ratios give strictly negative physical flux, contrary
to (3.1).  If no layer changes \(H\), it is a global invariant.

\subsection{One active coordinate}

Let \(A=n\) be active and let \(B,C\) remain in a finite box.  Every transition
is
\[
 (n,\phi)\to(n+k,\phi'),\qquad |k|\le2,
\]
with rate
\[
 c_2(\phi,e)n(n-1)+c_1(\phi,e)n+c_0(\phi,e).
\]
If \(2A\) is a source, every genuine reaction from it lowers \(A\), while
upward reactions have degree at most one:
\[
 \cL A\le-c_2A(A-1)+c_1A+c_0,\qquad c_2>0.
\]

Without \(2A\), every degree-one edge has increment zero or minus one.
Finite recurrent degree-one classes have strict negative mean and a Poisson
corrector, or are edgewise zero and are contracted.  A degree-zero birth
creates at most one \(A\) and an actual target \(A\) or \(A+D\); the creator
linkage provides its service path and the aggregate-debt theorem includes
every bounded-coordinate interruption.

In an all-zero class, let \(K\) be the bounded service species occurring in
one-\(A\) complexes.  Either a unary one-\(A\) complex or a zero-\(A\)
complex containing \(K\) gives unpaired service, or
\[
 A-\sum_{D\in K}D
\]
is an exact invariant.  Since the bounded coordinates and phase correction
are bounded, escape is impossible.


\subsubsection{Exact finite-phase correction}

Let \(Q_1\) be the finite generator formed by the coefficients of \(n\) in
the degree-one phase transitions.  For an irreducible recurrent class let
\(\pi\) be its stationary probability and \(g(\phi)\le0\) its conditional
mean \(A\)-increment.  If one negative edge is present, then
\(\bar g=\pi g<0\).  Solve
\[
 Q_1\chi=-(g-\bar g)
\]
with one normalization.  Since the phase is finite, \(\chi\) is bounded and
\[
 \cL\,[A+\chi(\phi)]=n\bar g+O(1),
\]
uniformly over the bounded-coordinate phase.

If \(\bar g=0\), nonpositivity forces every positive-flow degree-one edge to
have zero increment.  Contract the class.  Degree-zero events and their
complete fast relaxations form a finite macrochain.  A positive birth is
kept at its natural rate; its actual one-\(A\) target starts the faster
carrier process.  Repeated births are handled by aggregate debt.

In an edgewise-zero macroclass, the service-token dichotomy is exact.  If
there is no unary one-\(A\) complex and no zero-\(A\) complex contains a
service species from a binary one-\(A\) complex, then every one-\(A\)
complex carries exactly one such token and every zero-\(A\) complex carries
none.  Therefore \(A-\sum_{D\in K}D\) is preserved reaction by reaction.
Otherwise the finite phase contains an unpaired service or an exit.  Thus a
zero mean with nonzero variance is never classified as an invariant.

\subsection{Two active coordinates and the finite atlas}

We first record the interface between the stochastic trace and the finite
classification.  Call a linkage \emph{available} when it falls in one of the
available branches of the certified top-complex alternative.  In the
two-active chart, an available linkage has one of the following witnesses:
a top source containing two active particles, a unary top source, or an
active-plus-\(C\) target whose service cofactor is present in a lower complex.
The first two witnesses are enabled throughout the padded chart.  In the
third case the exact finite \(C\)-phase either restores the cofactor and
enables service, or leaves the box/support.  Actual-target return prefixes and
the carrier-race lemma therefore give strict service or exit.  Consequently,
a closed edgewise-zero terminal phase can contain only flat or shielded
linkages.  If both linkages represented in it are shielded, the pair atlas
below applies.

Let \(A,B\) be active and \(C\) bounded.  Up to exchanging \(A,B\), four
workloads cover every chamber and wall:
\[
 (1,1,0),\quad(2,3,0),\quad(1,2,0),\quad(1,3,0).
\]
Three independent implementations exhaust all assignments of the ten
complexes to linkage 1, linkage 2 or unused:
\[
\begin{array}{lr}
\text{workload assignments}&187{,}488\\
\text{shielded assignments}&446\\
\text{common affine invariants}&382\\
\text{deficiency-zero assignments}&60\\
\text{ordered service assignments}&4\\
\text{unclassified}&0.
\end{array}
\]

Available linkages are handled by Proposition~\ref{prop:layers}.  A closed
shielded phase is therefore one of:
\begin{enumerate}[label=(\roman*)]
\item a common affine invariant positive on \(A,B\);
\item a weakly reversible deficiency-zero network;
\item one of the two service architectures
\[
 \{C,2C\}\ \&\ \{0,A,2A,B+C\},
\]
\[
 \{0,C,2C\}\ \&\ \{A,2A,B+C\}.
\]
\end{enumerate}
The deficiency-zero branch has a summable product-Poisson stationary measure.
In the service branch \(W=B-C\); the mixed linkage preserves \(W\), while a
finite-shell \(C\)-trial gives
\[
 W_\tau\le W_0-1
\]
or a declared promotion/exit.

\section{Proof of the main theorem}

Assume, for contradiction, infinite expected positive return count.  Apply
Theorem~\ref{thm:green} and Lemma~\ref{lem:chart}.  The terminal chart has
nonnegative outward workload balance.

If it has three active coordinates, the strict source-layer flux is negative.
If it has two, Proposition~\ref{prop:layers} or the certified atlas gives
negative flux, an invariant, product-form stationarity, or service descent.
If it has one, the polynomial phase analysis gives quadratic descent, strict
corrected linear descent, an invariant, or exit.  With no active coordinates
the chart is finite.  Every possibility contradicts terminal escape.
Therefore the embedded chain has finite expected positive return count.

At every nonabsorbing population state, some enabled falling factorial is a
positive integer.  Hence \(\Lambda(x)\ge\kappa_{\min}>0\), and finite expected
embedded return count implies finite expected physical return time.

Population-increasing reactions have source molecularity at most one, so
their aggregate rate is bounded by \(C(1+|X|_1)\).  A linear pure-birth
comparison bounds upward jumps on finite time intervals.  Between upward
jumps total population is nonincreasing and the chain remains in a finite
shell, where total rate is bounded.  Thus the CTMC is nonexplosive.

The standard irreducible countable-state equivalence now gives positive
recurrence and uniqueness of the stationary probability
\cite{Norris1997}.  This proves Theorem~\ref{thm:main}. \qed

\appendix
\section{Forbidden activation regressions}

The proof never conditions on a future activation.  For
\[
 0\to2A\to A\to0
\]
the conditioned activation/path block satisfies
\[
 J_n=\frac{7\log n+1}{n^2}
      +O\!\left(\frac{\log n}{n^3}\right)>0.
\]
For
\[
 A\rightleftarrows2A,\qquad
 0\to A+B\to B\to0,
\]
conditioning on \(0\to A+B\) and appending only the target episode is
asymptotic to \(+\log n\).  Both invalid decompositions are preserved as
regression tests.  In the proof, every activation remains in its preceding
unconditioned physical source-layer trace.

\section{Verification boundary}

Computation is not used to prove the infinite-state Green or debt theorems.
It independently checks the finite atlas, priority graph alternatives,
one-active source signs, exact queue inequalities, Poisson arithmetic,
service supports, and the two activation counterexamples.


\begin{thebibliography}{9}
\bibitem{AndersonKim2018}
D.~F. Anderson and J.~Kim.
\newblock Some network conditions for positive recurrence of stochastically modeled reaction networks.
\newblock \emph{SIAM Journal on Applied Mathematics}, 78:2692--2713, 2018.

\bibitem{AndersonCappellettiKim2020}
D.~F. Anderson, D.~Cappelletti, and J.~Kim.
\newblock Stochastically modeled weakly reversible reaction networks with a single linkage class.
\newblock \emph{Journal of Applied Probability}, 57:792--810, 2020.

\bibitem{Kriebel2026Single}
A.~Kriebel.
\newblock Positive recurrence of bimolecular weakly reversible stochastic reaction networks with a single linkage class.
\newblock Version 0.2 manuscript, 2026.

\bibitem{Norris1997}
J.~R. Norris.
\newblock \emph{Markov Chains}.
\newblock Cambridge University Press, 1997.
\end{thebibliography}
\end{document}
